% Narrative companion to hash-bound export_report.py tables.
\section{Lossless storage and history engineering: October 30 update}
\label{sec:historyengineering}
\subsection{Scope and correctness}
The final engineering source is
\nolinkurl{d8eb51be69c5d6469ba12283f0cf811ba352345c}, local branch
\nolinkurl{evicastkv/storage-transport-20261030}. The sealed CPU suite records
2,631 passing tests, no failures/errors/skips. Both six-path GPU comparisons
complete all 12 workers and pass registered numerical/storage controls.
This is engineering readiness, not independent paper-quality confirmation.

All measurements reuse two consumed inputs: Description
\texttt{000000532947} and MM-Vet-route \texttt{v1\_197}, logical lengths
600/604. They use the pinned LLaVA-1.5 model in Section~\ref{sec:setup},
BF16, HF eager, batch one, one A100 per worker. The fixed 128-token schedule
uses the prefill prediction plus 127 decode forwards without early EOS stop.
Five warmups and 20 timed repetitions randomize paired format order within
each method/input worker. Intervals resample technical pairs, not independent
examples. No new benchmark identities or selection policies are introduced.

The six paths are an all-position lossless control, Prompt Top-K, local
append-only PrefixKV/VL-Cache adaptations, LOOK-M (R), and frozen
\path{mosaickv_exact}. LOOK-M retains its actual merged prompt K/V. The
all-position control retains every source position but encodes storage;
its memory is not ordinary uncompressed FullKV memory. Ordinary FullKV
remains the separate numerical/quality reference, excluded from latency.

\subsection{Mechanism: save a shared source, reconstruct exact K/V}
Source history stores encoded normalized layer inputs, logical maps, rotary
phases, and sparse bit corrections restoring original cached bits. A short
raw K/V tail and an additional normalized-input tail support online sealing.
Sealing commits transactionally only after original-K/V checks pass. Model
projections reconstruct unchanged values and positions. Sharing a projection
source saves history bytes without removing evidence, but costs recomputation.
This is a separately labeled lossless representation experiment, not prototype
merging or residual repair. The direct constructor removes an intermediate
format, reducing projection passes nine to four while retaining three final
original-K/V checks. The same optimization applies to the baseline paths.

\begin{table}[htbp]
\caption{Consumed-input storage in MiB. KV means ordinary lossless K/V
history within the same encoded-prompt path, not FullKV; source means
source-sealed history. Initial bytes are unchanged within each pair. Peak is
observed decode-owned storage, not whole-request HBM peak. These are format
differentials, not byte-matched selector comparisons.}
\label{tab:historystorage}
\resulttable{\linewidth}{history_storage}
\end{table}

For \Evi{}, final bytes fall 109.271 to 93.787~MiB on Description and
109.377 to 93.920~MiB on MM-Vet, approximately 14.2\%/14.1\%. Observed
decode-owned peak falls 119.474 to 103.990~MiB and 119.580 to 104.123~MiB,
approximately 13.0\%/12.9\%. Initial stays 61.543/61.688~MiB. This improves
final/peak storage within the same selection, not initial bytes or selection
quality.

\begin{figure}[htbp]
\centering
\begin{tikzpicture}
\begin{groupplot}[group style={group size=2 by 1,horizontal sep=1.2cm},
width=0.46\linewidth,height=5.0cm,ylabel={Owned storage (MiB)},
xtick={1,2,3},xticklabels={Initial,Final,Decode peak},ymin=0,ymax=130,
grid=major,legend style={font=\scriptsize,at={(0.03,0.97)},anchor=north west}]
\nextgroupplot[title={Consumed Description input}]
\addplot+[thick,color=prefixorange,mark=square*] table[x=boundary,y=ordinary] {generated/history_boundaries_llava_description.dat};
\addlegendentry{Ordinary lossless K/V history}
\addplot+[thick,color=eviblue,mark=*] table[x=boundary,y=source] {generated/history_boundaries_llava_description.dat};
\addlegendentry{Lossless source history}
\nextgroupplot[title={Consumed MM-Vet input},ylabel={}]
\addplot+[thick,color=prefixorange,mark=square*] table[x=boundary,y=ordinary] {generated/history_boundaries_mmvet_prefixkv.dat};
\addplot+[thick,color=eviblue,mark=*] table[x=boundary,y=source] {generated/history_boundaries_mmvet_prefixkv.dat};
\end{groupplot}
\end{tikzpicture}
\caption{EviCastKV format differential at unchanged selected evidence.
Lines identify accounting boundaries, not measurements over time.}
\label{fig:historyboundaries}
\end{figure}

The boundary charges payloads, source packets, corrections, phases, maps,
raw/normalized tails, layer leases, and observed constructor/reconstruction
tensors. It excludes inactive original prefill/captures, weights, attention
activations, internal GEMM workspace, and allocator reserve. Original prefill
can determine whole-request peak, which remains uncertified. Source-slot
capacity is fixed across paths but realized encoded bytes differ: this is
not a new common-physical-byte selection-policy comparison.

\subsection{Latency: constructor improvement within a slower storage path}
\begin{table}[htbp]
\caption{Fixed-schedule decode seconds. Source/KV comes from the K/V versus
source-history study; remaining columns come from a separate reference-source
versus direct-source constructor study. Do not combine medians across studies.
Intervals are ordinary technical-pair 95\% intervals, not population or
multiplicity-adjusted intervals.}
\label{tab:historylatency}
\resulttable{\linewidth}{history_latency}
\end{table}

Reference source history is 29.6--29.7\% slower than ordinary lossless K/V
history for \Evi{} here, roughly 29--36\% slower across compressed paths.
The direct constructor improves on that source-history reference: \Evi{}
p50 changes 11.742 to 11.282 seconds and 11.765 to 11.300 seconds, paired
reductions about 3.9--4.0\%, with identical initial/final/observed-peak bytes.
All ten within-path/input direct/reference upper confidence bounds are below
one. This does not prove no-slower execution against ordinary K/V history.

Timing includes online sealing/reconstruction but excludes original prefill,
selection, and initial prompt construction. Different methods use separate
GPU workers; cross-method medians are not paired latency wins. Earlier
roughly 12\% dual-projection decode and 57\% direct prompt-construction
improvements concern distinct references/phases; do not add their percentages
or subtract them from old end-to-end campaign times.

\subsection{What remains negative}
All 128 logits/tokens remain bitwise unchanged relative to each path's own
original selected/merged-cache transport, not necessarily FullKV. The
original 32-token model-reference outcomes remain different across selectors.

\begin{table}[htbp]
\caption{Consumed-input controls: R-L is model-generated FullKV-reference
ROUGE-L $\times100$, not annotation accuracy, corpus PPL, or standard MM-Vet
judge scoring. Bitwise preservation concerns 128 logits relative to each
path's original transport.}
\label{tab:historyquality}
\resulttable{0.95\linewidth}{history_quality_controls}
\end{table}

Description is a counterexample to universal superiority: the PrefixKV
adaptation has smaller final storage than \Evi{} with tied model-reference
quality. MM-Vet retains a quality deficit relative to FullKV for several
selections, including \Evi{}. The codec preserves, rather than repairs,
these outcomes. Negative quality interventions, short-history feasibility,
and the unpromoted checkpoint latency failure remain preserved. No fresh
campaign follows from the engineering pass alone.

\section{Paper-writing priorities}
\label{sec:writingactions}
\begin{enumerate}
\item Lead with the completed 256-image shared-byte COCO curves and separate
1,000/218 PrefixKV native/shared-byte tables. Retain all points and intervals,
Top-K ties/losses, and native final-memory/latency disadvantages. Do not use
the anomalous byte-rescaled adaptation as original-PrefixKV quality evidence.
\item Separate same-budget quality, equal-quality memory, initial savings,
and lifetime savings. A general ``30\% smaller than the nearest competitor
at preserved quality'' headline is not supported.
\item Put representation results in an engineering subsection/appendix with
the consumed identities, exact boundaries, lossless controls, and slower
ordinary-history comparison. Do not replace benchmark tables with diagnostics.
\item Keep claims limited to exact selection and separately labeled lossless
variants. Prototype/repair, other architectures, and native serving gain no
performance validation from this work.
\item Broader claims require a frozen candidate, certified memory boundary,
and disjoint paired native-metric validation, with common optimizations also
applied to competitors. These remain research requirements, not experiments
launched by this writing update.
\end{enumerate}
